\section{TỔNG QUAN TUẦN BÁO CÁO}\label{sec:tong_quan}

\subsection{Bối cảnh}

Đề tài mô phỏng một hệ thống mạng gồm Campus chính (Site 100), ba cơ sở Cần Thơ, Đà Nẵng, Nha Trang (Site 200, 300, 400), chi nhánh mới Site 500 và vùng điều khiển SD-WAN (Site 900) trên EVE-NG. Tại Site 100, bộ điều khiển Ryu quản lý sáu switch Open vSwitch bằng OpenFlow 1.3 qua VLAN quản trị 99. Các site liên kết với nhau bằng Cisco SD-WAN (Viptela 20.10.1) qua hai transport Internet và MPLS.

Kết thúc tuần trước, phần SDN đã có đủ số liệu năm mục tiêu và cơ chế fast-failover cho uplink Access; phía SD-WAN đã có công cụ đo độ tin cậy nhưng chưa có số liệu đo lặp. Tuần này nhóm thu số liệu SD-WAN đầu tiên, kiểm thử chức năng toàn hệ thống, thiết kế lại và đo failover cặp firewall HA ở Site 100, và bắt đầu một hạng mục mới: \textbf{Campus NetOps Bot} -- trợ lý vận hành mạng qua Telegram, chạy ngay trong topology và quản trị thiết bị theo kiểu in-band.

\subsection{Đề tài theo đề cương nhiệm vụ}

Đề tài số 2: \textbf{Xây dựng Campus Network sử dụng SDN và SD-WAN} (\textit{Implementation of an SDN- and SD-WAN-Based Campus Network}), ngành Mạng máy tính và truyền thông dữ liệu. Bảng~\ref{tab:de_cuong} đối chiếu năm nội dung của đề cương với phần đã hiện thực trong đồ án.

\begin{table}[H]
    \centering
    \small
    \caption{Nội dung đề cương và phần hiện thực trong đồ án}
    \label{tab:de_cuong}
    \begin{tabularx}{\textwidth}{|P{4.4cm}|X|P{2.0cm}|}
        \hline
        \textbf{Nội dung đề cương} & \textbf{Hiện thực} & \textbf{Trạng thái} \\
        \hline
        Mô hình Core -- Distribution -- Access & Site 100: 2 Core IOL (VRRP), 2 Distribution và 4 Access là Open vSwitch & Hoàn thành \\
        \hline
        VLAN theo khoa, phòng học, phòng máy, hành chính & VLAN 10 CNTT, 20 TTK, 30 Luật, 40 Hành chính, 90 Server Farm; mỗi chi nhánh có VLAN riêng & Hoàn thành \\
        \hline
        SDN Controller quản lý switch & Ryu + OpenFlow 1.3 điều khiển 6 OVS, tự phục hồi, fast-failover & Hoàn thành \\
        \hline
        SD-WAN nối Campus chính với cơ sở phụ & Cisco Viptela: 4 chi nhánh, 2 đường Internet và MPLS, BFD/AAR & Hoàn thành \\
        \hline
        STP/RSTP, OSPF, DHCP, ACL, dự phòng liên kết & Rapid-PVST trên Core Site 100, STP chi nhánh, OSPF lõi, DHCP tập trung, ACL trên firewall; Port-channel, VRRP, firewall HA & Hoàn thành \\
        \hline
    \end{tabularx}
\end{table}

Đề cương đặt ra bốn mục tiêu: (1) đơn giản hóa quá trình quản trị mạng Campus; (2) quản lý tập trung luồng dữ liệu giữa các khu vực; (3) tăng khả năng mở rộng khi thêm khoa hoặc cơ sở mới; (4) đảm bảo tính sẵn sàng cao cho các dịch vụ đào tạo. Bảng~\ref{tab:tieu_chi} cho biết nhóm đã có số liệu tới đâu với từng tiêu chí đánh giá; các tiêu chí chưa đủ số liệu là trọng tâm của kế hoạch tuần tới (Chương~\ref{sec:ke_hoach}).

\begin{table}[H]
    \centering
    \small
    \caption{Tiêu chí đánh giá theo đề cương và số liệu hiện có}
    \label{tab:tieu_chi}
    \begin{tabularx}{\textwidth}{|P{4.4cm}|X|P{2.0cm}|}
        \hline
        \textbf{Tiêu chí} & \textbf{Số liệu hiện có} & \textbf{Trạng thái} \\
        \hline
        Thời gian thêm mới VLAN hoặc khu vực mạng & Cập nhật 07/11: thêm VLAN cấu hình xong 5,3~s, máy trạm dùng được khoảng 40~s (IP tĩnh); chưa bấm giờ thêm site & Một phần \\
        \hline
        Thời gian khôi phục khi mất liên kết & SDN dưới 0,3~s; SD-WAN khoảng 8~s; firewall HA 11--21~s & Một phần \\
        \hline
        Hiệu năng giữa các VLAN và các cơ sở & Delay, jitter, mất gói 5 site. Cập nhật 06/11: thông lượng giữa vEdge khoảng 2~Mbit/s (giới hạn WAN mô phỏng); chưa đo giữa máy trạm & Một phần \\
        \hline
        Tải trên Core và Distribution & Chưa đo & Chưa đo \\
        \hline
        Khả năng quản lý chính sách tập trung & vSmart, Ryu, NetOps Bot; chưa đo định lượng & Một phần \\
        \hline
    \end{tabularx}
\end{table}

\subsection{Mục tiêu đặt ra cho tuần}

Theo kế hoạch đã trình bày ở báo cáo tuần trước:

\begin{enumerate}
    \item Đo lặp SD-WAN: failover link WAN, thời gian khôi phục khi mất vSmart, độ sẵn sàng.
    \item Kiểm thử failover cặp firewall ASAv HA và VRRP trên Core.
    \item Bắt đầu hoàn thiện Site 500 (vEdge, switch, máy trạm, chính sách).
    \item SDN: thử fast-failover cho uplink Distribution.
    \item Theo dõi và thu bằng chứng nếu vEdge mất kết nối.
    \item Đưa số liệu vào bản nháp báo cáo cuối kỳ.
\end{enumerate}

Trong tuần, nhóm bổ sung mục tiêu mới: xây dựng công cụ vận hành tự động (NetOps Bot) để phục vụ chẩn đoán lỗi và phần demo khi bảo vệ.

\subsection{Đối chiếu kế hoạch và thực tế}

\begin{table}[H]
    \centering
    \small
    \caption{Đối chiếu kế hoạch tuần và kết quả thực tế}
    \label{tab:doi_chieu}
    \begin{tabularx}{\textwidth}{|c|X|P{2.4cm}|P{4.6cm}|}
        \hline
        \textbf{\#} & \textbf{Đầu việc theo kế hoạch} & \textbf{Trạng thái} & \textbf{Ghi chú} \\
        \hline
        1 & Đo lặp SD-WAN & Hoàn thành một phần & Có RTO/PRR cho vEdge 1 Site 200 (2 lần), so sánh hiệu năng Campus chính -- chi nhánh, bắt đầu đo độ sẵn sàng; chưa đủ 3--5 lần mỗi kịch bản \\
        \hline
        2 & Failover firewall HA và VRRP & Hoàn thành & Thiết kế lại HA Site 100 (transit VLAN, SW-WAN-Edge1/2, OSPF với vEdge), đo 4 kịch bản; sửa VRRP split-brain \\
        \hline
        3 & Hoàn thiện Site 500 & Hoàn thành & 2 vEdge Site 500 reachable trên vManage; Brand-FW-S500 (VLAN 10/99, DHCP), SW-S500, máy trạm nhận IP; sửa lỗi Site 500 làm mất Internet các site \\
        \hline
        4 & Fast-failover uplink Distribution & Chưa thực hiện & Ưu tiên firewall HA; đưa vào phần đo thời gian khôi phục tuần tới \\
        \hline
        5 & Theo dõi vEdge mất kết nối & Hoàn thành & 12/12 thiết bị SD-WAN reachable, BFD 92/92; NetOps Bot nay tự cảnh báo khi vEdge mất kết nối \\
        \hline
        6 & Bản nháp báo cáo cuối kỳ & Chưa thực hiện & Bảng tiến độ và README đã cập nhật \\
        \hline
        7 & \textit{(Phát sinh)} Campus NetOps Bot & Hoàn thành & Chạy thật trên Telegram, quản trị in-band Site 100 và 3 chi nhánh, selftest 23/23 \\
        \hline
    \end{tabularx}
\end{table}

Nhóm hoàn thành ba trên sáu đầu việc (firewall HA, Site 500, theo dõi vEdge), làm được một phần việc đo lặp SD-WAN, và phát sinh thêm NetOps Bot cùng kiểm thử chức năng toàn hệ thống. Fast-failover Distribution và bản nháp báo cáo cuối kỳ chưa làm; tuần tới nhóm tập trung đo các chỉ số theo tiêu chí đề cương (Chương~\ref{sec:ke_hoach}).
